\documentclass[11pt]{article}
\usepackage[a4paper,margin=1in]{geometry}
\usepackage{amsmath,amssymb,mathtools,bm}
\usepackage{enumitem,booktabs,array}
\usepackage{tcolorbox}
\usepackage{hyperref}
\usepackage[protrusion=true,expansion=false]{microtype}
\hypersetup{colorlinks=true,linkcolor=blue,urlcolor=blue,citecolor=blue}
\setlist[itemize]{topsep=2pt,itemsep=2pt}
\setlist[enumerate]{topsep=2pt,itemsep=2pt}
\newtcolorbox{keybox}{colback=blue!3!white,colframe=blue!40!black,boxrule=0.4pt,arc=1mm}
\newtcolorbox{alertbox}{colback=red!3!white,colframe=red!40!black,boxrule=0.4pt,arc=1mm}
\newtcolorbox{answerbox}{colback=green!3!white,colframe=green!40!black,boxrule=0.4pt,arc=1mm}
\newtcolorbox{drillbox}{colback=yellow!5!white,colframe=orange!60!black,boxrule=0.4pt,arc=1mm}
\newcommand{\EV}{\mathbb{E}}
\newcommand{\Var}{\mathrm{Var}}
\newcommand{\blank}[1]{\underline{\hspace{#1}}}

\title{W01-01: Kashyap, Rajan, Stein (2002, JF) \\
\large ``Banks as Liquidity Providers: An Explanation for the Coexistence of Lending and Deposit-Taking'' --- Active Recall Workbook}
\author{Banking \& Risk Management --- Exam Prep}
\date{\today}
\begin{document}\maketitle

\begin{keybox}
\textbf{Paper in one sentence.} Deposits and loan-commitment (credit-line) drawdowns are \emph{imperfectly correlated}, so a single pool of liquid assets can back both: that \emph{synergy} is why lending and deposit-taking co-exist inside the same institution.

\textbf{Placement.} Foundational paper for Part~I of the course: provides the modern theoretical rationale for the bank as a combined deposit-and-loan franchise. Direct empirical sequel is Gatev--Strahan (W01-02).
\end{keybox}

\section*{Round 1: Complete Walk-Through}

\subsection*{1.1 Core insight}
Historically banks were explained either as \emph{liquidity providers to depositors} (Diamond--Dybvig) or as \emph{monitors of borrowers} (Diamond 1984). KRS ask: why the same firm? Their answer: a bank makes liquidity promises on \emph{both sides} of its balance sheet --- to depositors who may withdraw and to firms that hold undrawn credit lines. If withdrawals $W$ and drawdowns $D$ are \emph{less than perfectly correlated}, the liquidity reserve needed to back both is less than the sum of what two separate institutions would need.

\subsection*{1.2 Model sketch}
Let a bank hold liquid asset $L$, face stochastic deposit outflows $\tilde W$ and commitment drawdowns $\tilde D$. The bank is solvent if
\[
L \ge \tilde W + \tilde D.
\]
If $\rho = \text{Corr}(\tilde W, \tilde D) < 1$, then for a given failure probability the required $L^*$ satisfies
\[
L^* \;<\; L^*_W + L^*_D,
\]
because $\Var(\tilde W + \tilde D) = \sigma_W^2 + \sigma_D^2 + 2\rho\sigma_W\sigma_D < \sigma_W^2 + \sigma_D^2$ when $\rho<0$ or small. The synergy is economically meaningful especially when $\rho$ is \emph{negative} --- i.e., when deposits flow \emph{in} at the same time firms draw down lines.

\subsection*{1.3 Empirical content}
\begin{itemize}
\item Cross-section of US banks 1992--1996: deposit-to-commitment ratio is strongly positive across banks, and banks with more transaction deposits also have more unused commitments.
\item During periods of tight liquidity (e.g., 1998 LTCM episode hinted at but Gatev--Strahan nails this), deposit inflows and commitment drawdowns move in \emph{opposite} directions $\Rightarrow$ the synergy is maximal exactly when it is most needed.
\item Banks with transaction deposits (checking accounts) hold less liquid securities relative to commitments than banks with time deposits --- consistent with the hedging story.
\end{itemize}

\subsection*{1.4 Identification / econometric considerations}
\begin{alertbox}
\begin{itemize}
\item Endogeneity: bank capital structure and commitment volume are jointly chosen. Cross-sectional correlations are suggestive, not causal.
\item The model predicts a \emph{diversification benefit} only when $\rho(\tilde W,\tilde D)$ is low. Would fail in a systemic-run state ($\rho\to 1$) --- consistent with why 2008 crisis forced government backstops.
\item Deposit insurance and the discount window make deposits \emph{sticky}, which raises the correlation of ``effective'' liquidity shocks. The synergy is largest in insured banks.
\end{itemize}
\end{alertbox}

\section*{Round 2: Level 1 Fill-in}
\begin{enumerate}
\item A bank is solvent iff $L \ge \blank{1cm} + \blank{1cm}$.
\item $\Var(\tilde W + \tilde D) = \sigma_W^2 + \sigma_D^2 + \blank{2cm}$.
\item Synergy is largest when $\rho(\tilde W,\tilde D)$ is \blank{1.5cm}.
\item Banks with more \blank{1.5cm} deposits have a stronger deposit-commitment hedge.
\item Period when deposits flow in while lines are drawn: \blank{2cm}.
\end{enumerate}

\section*{Round 3: Level 2 Prompts}
\begin{enumerate}
\item Derive the required liquid-asset buffer $L^*(\alpha)$ such that $\Pr(L<\tilde W+\tilde D)=\alpha$, assuming $(\tilde W,\tilde D)$ joint-normal. Show that $\partial L^*/\partial\rho > 0$.
\item Explain why Diamond--Dybvig alone cannot justify deposit-plus-lending inside a single bank.
\item What breaks the KRS synergy during a systemic run? Draw the implications for optimal regulatory capital.
\item Relate KRS to the later Gatev--Strahan (2006) empirical finding.
\end{enumerate}

\section*{Round 4: Minimal Scaffolding}
From a blank page, reproduce: (i) the core synergy argument, (ii) the variance-of-sum decomposition, (iii) three empirical predictions, (iv) two regulatory implications.

\section*{Round 5: Blank Page}
Why is the deposit-plus-loan-commitment bank the \emph{natural} institutional form in a world where liquidity shocks are imperfectly correlated? Write a one-page essay.

\vspace{6pt}
\begin{drillbox}
\textbf{Speed Drill.} (1) Variance of sum formula. (2) Sign of $\rho$ needed for synergy. (3) Why transaction deposits are the right hedge. (4) What happens to the synergy in a systemic run? (5) Diamond--Dybvig vs.\ KRS rationale in one sentence each.
\end{drillbox}

\section*{Quick Reference \& Answer Key}
\begin{answerbox}
\begin{itemize}
\item \textbf{Core identity.} $\Var(\tilde W + \tilde D) = \sigma_W^2 + \sigma_D^2 + 2\rho\sigma_W\sigma_D$.
\item \textbf{Synergy.} $L^*(\rho)<L^*_W+L^*_D$ whenever $\rho<1$; strongest for $\rho<0$.
\item \textbf{Empirical support.} Banks with more demand/transaction deposits have \emph{more} unused loan commitments; in flights-to-quality, deposits rise while lines are drawn.
\item \textbf{KRS vs.\ alternatives.} Diamond--Dybvig (1983) explains deposits only; Diamond (1984) explains loan-monitoring only. KRS is the joint explanation.
\end{itemize}
\end{answerbox}

\begin{keybox}
\textbf{30-second exam pitch.} KRS (2002) is the theoretical anchor for treating banks as combined liquidity providers on both sides of the balance sheet. Deposits can be withdrawn on demand; loan commitments can be drawn on demand. Because the two liquidity shocks are imperfectly --- and often \emph{negatively} --- correlated, a single liquid-asset buffer backs both at lower cost than two separate institutions would require. The variance-of-sum identity $\Var(\tilde W+\tilde D)=\sigma_W^2+\sigma_D^2+2\rho\sigma_W\sigma_D$ makes the mechanism transparent. Empirically, transaction-deposit-heavy banks hold relatively more undrawn commitments, and during flights-to-quality deposits flow in while credit lines are drawn --- exactly the negative-$\rho$ state in which the synergy is largest. The synergy breaks in a systemic run where $\rho\to 1$, which is why the backstop (deposit insurance, discount window, LOLR) is integral to the combined institution.
\end{keybox}

\end{document}
